\section{Equality orbits instead of a finite test alphabet}
\label{sec:orbits}

\subsection{The distinction that finite tests miss}
A cache can have only two storage slots and still process arbitrarily many
different keys. Its finite memory does not imply a finite concrete state space:
the slots may contain any two atoms. Testing all request sequences over two
chosen keys also does not represent all behaviors. It removes the possibility
of receiving a third key while both slots are occupied.

The appropriate finite state is an \emph{equality pattern}, including the
relationships between both implementations being compared. Renamings of atom
identities preserve execution because the allowed operations cannot inspect
anything except equality. The general perspective belongs to nominal automata
and permutation symmetries~\cite{nominal}. This section gives a direct,
restricted construction whose entire search and replay are included.

\subsection{A small, typed machine language}
Fix an infinite set $\Atoms$ of atoms, a finite tag set $T$, and distinct fixed
constants $c_0,\ldots,c_{c-1}\in\Atoms$. A machine has a finite control $q$ and
$k$ optional registers, each containing either an atom or $\bot$ (empty).
An input is $(t,u)\in T\times\Atoms$; $\bot$ is \emph{not} an input atom.

A selector is one of: current input, a specified old register, a specified
fixed constant, or empty. Guards are Boolean combinations of equality between
selectors. Register updates simultaneously assign one selector to each register.
An output is either a Boolean guard value or an optional atom selected from the
same old-state/input environment. The output constructors are tagged, so
Boolean false and atom zero cannot accidentally compare equal.

At each control and tag, the finite rule list is ordered. The first enabled
rule is chosen. The admitted format requires an unconditional default rule for
each control/tag cell; this makes execution total by construction. Rules after
a default in the same cell are simply unreachable under this semantics. A
consumer must not reorder rules or reinterpret the list as nondeterministic.
Initial registers contain only constants or empty. Extra fixed initial atoms
can be represented by adding constants, provided their distinctness and their
source interpretation are proved.

The executable representation uses nonnegative integers as \emph{names} for
atoms, but no arithmetic or order operation on those integers is in the machine
language. A Python sorting operation used to choose representatives is not an
operation available to the modeled program.

\subsection{Joint support and canonical representatives}
For machines $A$ and $B$ with $k_A$ and $k_B$ registers, form the paired state
\[
 s=(q_A,q_B,r_A,r_B).
\]
Let $k=k_A+k_B$. Its support consists of the distinct nonconstant atoms in the
concatenated tuple $(r_A,r_B)$; the fixed constants are retained separately.
Two paired states are equivalent when their controls match and a bijection of
atoms fixing every constant carries one \emph{joint} register tuple to the
other. Empty is also fixed.

For example, with no constants,
\[
 ((a,b),(b,a))\quad\text{and}\quad((17,6),(6,17))
\]
represent the same paired orbit. The pairs $((a),(a))$ and $((a),(b))$ with
$a\ne b$ do not. Canonicalizing each machine's single register separately
would map both pairs to $((0),(0))$, losing the exact fact that their outputs
might disagree. The canonicalizer must see both states at once.

The producer preserves the integer labels of constants and scans the joint
tuple left-to-right, assigning consecutive fresh labels to previously unseen
nonconstant atoms. Empty remains empty. Canonicalization is a search
optimization; the checker uses a different representation: the full pairwise
equality matrix of
\[
 (\bot,c_0,\ldots,c_{c-1},r_{A,0},\ldots,r_{B,k_B-1}).
\]
It therefore accepts noncanonical labels while preserving all cross-machine
relations and all anchors.

\begin{lemma}[Equality patterns characterize orbits]
\label{lem:pattern}
Two optional-register tuples have the same equality matrix with the same fixed
anchors if and only if a permutation of $\Atoms$ fixing those anchors carries
one tuple to the other.
\end{lemma}
\begin{proof}
A permutation preserves equality, constants, and empty, proving one direction.
In the other direction, equality of matrices defines a well-defined bijection
between the finitely many nonconstant atoms appearing in the tuples. It maps
each left occurrence to the corresponding right occurrence. A bijection between
two finite subsets extends to a permutation of their union, and hence to a
permutation of all atoms by fixing the complement. Constants remain fixed
because their equality classes were anchored separately.
\end{proof}

\subsection{Why one fresh representative is enough}
Given a paired state, enumerate as possible input atoms each live atom in either
machine, each fixed constant, and one atom distinct from all of them. Tags are
enumerated independently.

\begin{lemma}[Fresh-input coverage]
\label{lem:freshinput}
For a fixed paired state, every input atom belongs to exactly one of the
following relevant cases: a particular fixed/live atom, or the single class of
atoms outside the fixed/live support. Every atom in the last class produces a
successor and output pair related by a permutation fixing the current state.
\end{lemma}
\begin{proof}
Two atoms outside the support can be exchanged while fixing the support and
constants. By induction on selector and guard syntax, evaluation commutes with
such permutations: equality is preserved, Boolean connectives are unchanged,
and a selector either returns an anchored old value or the renamed input.
The first enabled rule is therefore unchanged. Simultaneous copying and
selected outputs commute with the permutation as well. Hence one chosen fresh
atom represents all outside-support atoms.
\end{proof}

``Fresh'' here means absent from \emph{currently remembered} registers and fixed
constants, not never used at any earlier point in the execution. The machine
cannot distinguish an old forgotten name from a globally new name. Insisting
on globally new representatives would enlarge witness alphabets unnecessarily;
merging a live name into the fresh case would be unsound.

For atom-valued outputs, compare the two actual selected values under the same
representative input \emph{before} forgetting their names. The property being
proved is equality of the two outputs at each time, not a separate renaming of
each output stream. No history of previous outputs is required for pointwise
equality. A specification that tests equality with an arbitrarily old output
would need to represent that additional memory explicitly.

\subsection{Search and finite certificates}
The producer performs breadth-first search from the paired initial orbit.
For every stored representative, every tag, and every representative atom,
it computes both outputs and both successors. Different outputs yield a
counterexample. Otherwise it inserts the successor's \emph{joint} canonical
state. A parent pointer records the input case and predecessor for each newly
visited orbit.

\begin{lstlisting}
seen := { joint_initial_orbit }
pending := FIFO queue containing it
while pending is not empty:
    s := pop_front(pending)
    for tag in all_tags:
        for a in constants_and_live_atoms(s) + [one_fresh_atom(s)]:
            (outA, nextA) := execute(A, left(s), tag, a)
            (outB, nextB) := execute(B, right(s), tag, a)
            if outA != outB:
                lift_parent_path_to_concrete_inputs(s, tag, a)
                return Different(word)
            t := joint_canonical(nextA, nextB)
            if t not in seen:
                record its parent; add t to seen and pending
return Equivalent(list_of_seen_representatives)
\end{lstlisting}

A positive receipt is simply a finite list $S$ of paired orbit representatives.
The checker independently verifies that the initial equality pattern belongs
to $S$, every representative input produces equal outputs, and every successor
pattern belongs to $S$. It does not trust producer-supplied edge lists, a
canonicalization hash, a minimality claim, or a search status.

\begin{theorem}[Orbit receipt soundness]
\label{thm:orbitsound}
An accepted positive receipt proves that $A$ and $B$ emit identical output
sequences for every finite input word over $T\times\Atoms$.
\end{theorem}
\begin{proof}
Relate concrete paired states to the listed representative with the same
controls and anchored equality matrix. The initial pair is related. For a
concrete input, use Lemmas~\ref{lem:pattern} and~\ref{lem:freshinput} to
transport the pair and input to a checked representative case. That case has
equal outputs and a successor represented in $S$. Transporting back preserves
output equality and the successor relation. Induction on word length proves
pointwise equality of the entire output sequences.
\end{proof}

A negative receipt is a word of ordinary tags and concrete nonnegative atom
labels. The independent checker executes the original two machines and finds
an output mismatch. During witness construction the producer lifts each
abstract input relative to the \emph{current concrete joint state}; it does not
concatenate canonical names from different states as though they were globally
consistent identities. The latter shortcut is a common incorrect use of
quotient paths.

\subsection{Termination, cutoff, and witness length}
Let $B_n$ denote the $n$th Bell number, the number of partitions of $n$ labeled
positions. A tuple of $k$ optional registers with $c$ fixed constants is
determined by a partition of $k+c+1$ positions, with the anchor positions kept
pairwise distinct. Thus $B_{k+c+1}$ is a simple upper bound; it deliberately
counts some impossible anchor identifications. With $c=0$, the empty anchor
makes $B_{k+1}$ the exact number of optional equality patterns available over an
infinite atom set.

\begin{theorem}[Finite orbit decision]
\label{thm:orbitcomplete}
Without cutoffs, the breadth-first search terminates after visiting at most
\[
 N=|Q_A|\,|Q_B|\,B_{k_A+k_B+c+1}
\]
paired orbits. It either returns a positive receipt or a shortest differing
input word. If $R$ paired orbits are reachable, a differing word, when one
exists, has length at most $R$.
\end{theorem}
\begin{proof}
There are at most $N$ control-pattern pairs, each with finitely many input
cases. If all are exhausted without a mismatch, the visited set satisfies
Theorem~\ref{thm:orbitsound}. If a mismatch is found, path lifting gives a
concrete differing word. Conversely every concrete path projects to this orbit
graph, so a concrete mismatch cannot be missed.

An output mismatch occurs on an outgoing edge. A shortest path to the source
of such an edge has no repeated paired orbit: a repeated segment can be removed,
and the remaining abstract path can still be lifted. The prefix has length at
most $R-1$, giving total length at most $R$. Breadth-first traversal considers
all shorter prefixes first, so its first mismatch has minimum word length.
\end{proof}

This is not a claim that the implementation visits every orbit up to the
Bell-number bound. Often almost all are unreachable. The cache-renaming
example needs only three paired orbits despite the crude bound $B_5=52$.

\begin{corollary}[A justified finite atom carrier]
\label{cor:cutoff}
Let $F\subseteq\Atoms$ contain all fixed constants and have size at least
$k_A+k_B+c+1$. The two admitted machines are equivalent over $\Atoms$ if and
only if they are equivalent on all words over $T\times F$.
\end{corollary}
\begin{proof}
One implication is restriction. For the other, lift any differing orbit path
using names in $F$. At each point at most $k_A+k_B$ nonconstant names are live,
so $F$ contains a name outside the live/fixed support. Use that name for the
fresh case, reusing names only after they have been forgotten. This preserves
every required equality and inequality along the path and produces a differing
word over $F$.
\end{proof}

The result is a completeness theorem for \emph{all words over a derived
carrier}, not for all words up to a chosen length. It also assumes the specified
copy/equality language and anchored initial registers. It does not justify a
similar cutoff for arbitrary programs that compare keys by hash or order,
remember an unbounded set, or create fresh atoms internally with additional
semantics.

\subsection{Worked example: the cache distinction}
The fixtures compare a two-slot least-recently-used cache with a first-in,
first-out cache. Both start empty, both report whether an input key was present
before the request, and both insert a missed key at the front while dropping
the old tail. On a hit in the tail, LRU moves the hit to the front while FIFO
leaves the insertion order unchanged. Head hits leave both unchanged.

The search returns the word
\[
 0,\ 1,\ 0,\ 2,\ 0.
\]
The first two requests fill the cache. The third makes 0 most recent under LRU,
but does not change FIFO's insertion order. Request 2 then evicts 1 under LRU
and 0 under FIFO. The final request is a hit only under LRU. The delivered
search reports five visited pair orbits and eleven checked representative
edges before finding the mismatch.

\begin{center}
\begin{tabular}{@{}rccccc@{}}
\toprule
Request & 0 & 1 & 0 & 2 & 0 \\
\midrule
LRU hit output & false & false & true & false & true \\
FIFO hit output & false & false & true & false & false \\
\bottomrule
\end{tabular}
\end{center}

Exhaustive concrete exploration over \emph{only two keys} visits seven paired
states and reports equivalence. That answer is correct for that smaller
alphabet: after both keys have appeared, all later requests hit. It is wrong
as a conclusion about arbitrary keys. The fresh representative is the missing
semantic case, not merely a request for a larger random sample.

As a positive control, swapping the two physical registers and consistently
rewriting all selector indices produces an equivalent implementation. Its
certificate has three joint representatives:
\[
 (\bot,\bot\mid\bot,\bot),\quad
 (0,\bot\mid\bot,0),\quad
 (0,1\mid1,0).
\]
There are six representative input checks. The same three-state relation also
certifies a variant that returns the evicted atom rather than a hit Boolean.
This exercises cross-machine equality, empty registers, and atom-valued
outputs, not just control-state isomorphism.

\subsection{Checking cost and representation choices}
Let $L$ bound the cost of interpreting the rule list and expressions for one
transition, $k=k_A+k_B$, and $R$ be the receipt's orbit count. Each orbit has at
most $c+k+1$ representative atoms per tag. With hashed equality signatures,
the checker costs roughly
\[
 O\bigl(R\,|T|\,(c+k+1)\,(L+(c+k)^2)\bigr)
\]
primitive operations, plus input validation. The quadratic term is the simple
pairwise equality-matrix construction, intentionally distinct from the
producer's linear scan canonicalizer. A verified partition encoding could
reduce it later. Bell-number growth remains the limiting factor as register
count rises; the construction is not advertised as a general scalable
replacement for every symbolic automata algorithm.
